\section{Logika Boolean dalam Relung Arsitektur Perangkat Lunak}

Beranjak dari konteks filosofis, bagian ini membahas penerapan logika langsung di layar terminal komputasi \cite{levin_discrete,saylor_cs202}. Pengembangan perangkat lunak di abad ke-21 bergantung pada \textbf{Aljabar Logika Boolean}.

\subsection{Boolean sebagai Dasar \textit{Control Flow}}

Bila jantung seorang manusia dialiri plasma sel darah, maka struktur \textit{Control Flow} dan \textit{Branching Decision} (Percabangan Keputusan) dari setiap program di bumi ini dialiri murni oleh denyut nadi \textit{Boolean} (Nilai limit biner 1 untuk T, 0 untuk F). 

Tanpa evaluasi logika, skrip perangkat lunak hanya berjalan linier dari baris 1, 2, 3 hingga selesai, tanpa percabangan. Evaluasi syarat logika (\texttt{IF}, \texttt{ELSE IF}, repetisi \texttt{WHILE}) memungkinkan percabangan alur eksekusi. Mesin kini dapat mengelola aliran berdasarkan hasil \texttt{True} dan \texttt{False} \cite{levin_discrete}.

\begin{figure}[htbp]
\centering
\begin{tikzpicture}[
  node distance=0.8cm,
  decision/.style={diamond, draw=black, fill=red!15, aspect=2, align=center, font=\bfseries\small, text width=2.2cm},
  process/.style={rectangle, draw=black, fill=blue!15, rounded corners, minimum height=0.8cm, text width=2.8cm, align=center},
  start/.style={ellipse, draw=black, fill=green!20, font=\bfseries},
  arrow/.style={->, thick, >=stealth}
]
  \node[start] (start) {Mesin Aktif Start};
  \node[decision, below=of start] (cond) {Verifikasi Operator Boolean Bersyarat?};
  \node[process, below left=1cm and 1.2cm of cond] (t) {Eksekusi Jalur Terang (Blok True)};
  \node[process, below right=1cm and 1.2cm of cond] (f) {Bypass / Buang (Blok False)};
  \node[process, below=2.5cm of cond] (end) {Penyatuan Muara Akhir Lanjut};
  \draw[arrow] (start) -- (cond);
  \draw[arrow] (cond) -| node[near start, left, font=\bfseries, text=blue] {Ya Sesuai!} (t);
  \draw[arrow] (cond) -| node[near start, right, font=\bfseries, text=red] {Palsu / Tidak!} (f);
  \draw[arrow] (t) |- (end);
  \draw[arrow] (f) |- (end);
\end{tikzpicture}
\caption{Visualisasi organ gerbang \textit{Cyber} evaluasi Boolean memotong jalan eksekusi.}
\label{fig:alur-boolean}
\end{figure}

\subsection{Merajut Kondisi Kompleks Dengan Senjata Operator}

Meminta mesin mendeteksi kualifikasi pelamar kerja yang gajinya cuma di bawah \textit{UMR} adalah kasus sepele. Di industri \textit{Google} dan \textit{Mobile Banking}, algoritma pelindung menyortir puluhan filter maut secara bersamaan menuntut ikatan \textbf{Rantai Operator \texttt{AND (Konjungsi), OR (Disjungsi), NOT (Negasi)}} \cite{python_docs_boolean}.

\begin{contoh}[Simulasi Skrip Realita Mesin]
Kasus \textit{Cyber Security}: Akses Brankas Transfer Dana diloloskan hanya bagi "Pengguna ber-IP domestik (\textbf{ATAU} berstatus staf VIP)" \textbf{SERTA mutlak jua} "Usianya TIDAK bergolong Balita Bebas Rekening".

\begin{lstlisting}[style=pythonstyle, caption={Operasi Filter Keamanan dengan Operator Boolean}]
# Skema Filter Keamanan Finansial:
staf_vip = False
ip_domestik = True
usia_akun_tahun = 3

# Merakit rute majemuk evaluasi Boolean mutlak dewa:
# Status evaluasi -> (True OR False) AND (NOT False)
# Status evaluasi -> (True) AND (True)
izin_lolos = (ip_domestik or staf_vip) and not (usia_akun_tahun < 12)

if izin_lolos:
    print("GERBANG BRANKAS PLATINUM DIBUKA! Validasi Sah.")
else:
    print("PENYUSUP BLOKIR! Anda terlempar ke Penjara Alarm.")
\end{lstlisting}
Melalui penguasaan hukum ekuivalensi (seperti Hukum De Morgan), insinyur basis data dapat menyederhanakan ekspresi \texttt{if} kompleks bermuatan puluhan baris menjadi beberapa baris yang efisien.
\end{contoh}
